\section{Experiments}
\label{sec:setup}

\subsection{Benchmarks}
\label{sec:benchmarks}

\paragraph{UMR} Uniform Meaning Representation \citep{vangysel2021designing,bonn2024building} combines sentence-level concept and argument graphs, word alignments, and document-level temporal, modal, and coreference relations. Its guideline has 365{,}742 characters. Resources include Chinese and English PropBank frames \citep{xue2009adding,palmer2005propbank}, lists of rolesets, relations and attribute values, and the guideline itself. We evaluate Chinese and English annotation using the first UMR parsing shared task \citep{stepanek2026first}.

\paragraph{TimeML} TimeML \citep{pustejovsky2003timeml} requires recognizing events and time expressions, assigning their attributes and normalized values, and linking them with fourteen temporal relation types within and across sentences. We follow the 71-page TimeML 1.2.1 guideline \citep{sauri2006timeml}, which has 115{,}146 characters and is the only resource file, and evaluate on TempEval-3 \citep{uzzaman2013semeval}. Outputs use inline entity tags followed by temporal links. Data selection and training/test splits for both benchmarks are detailed in Appendix~\ref{app:data}.

\subsection{Judges and evaluation}
\label{sec:evaluation}

\paragraph{UMR} We use AnCast++ \citep{sun2024anchor,sun2025ancastpp} to align predicted and gold nodes, comparing concepts, senses, attributes, arguments, and document-level relations (Appendix~\ref{app:judge}). We report AnCast++ sentence-graph F (\emph{F}, the main metric) and its comprehensive score including document-level relations (\emph{Comp.}), after the cleanup of \citet{chun2024pipeline}, macro-averaged over documents. We also report the shared task's official \jumach{} score \citep{stepanek2026first,zeman2026word}, computed with \texttt{compare\_umr.pl} against the official test gold. It maps nodes through token alignment and therefore also scores alignment; its gold differs from the cleaned gold used by AnCast++ (Appendix~\ref{app:data}).

\paragraph{TimeML} We built timemljudge using TempEval-3's extent and attribute measures for events and times \citep{uzzaman2013semeval} and by extending TempEval-2's categories of temporal links \citep{verhagen2010semeval} to event--time and event--event within or across sentences, event--document time, and time--time. Each category has separate checks for which pairs are linked and whether their relations are correct. For overall evaluation we use the official TempEval-3 scorer, reporting temporal-awareness TLINK F1 (\emph{TLINK}, the main metric), event and time-expression extent F1, and time-expression value F1.

\subsection{Setup}

\paragraph{Models}
The annotators are Claude Haiku~5.5 and GPT-6.1-sol. Claude Sonnet~5.5 writes the seed and is the reflector and editor in every round. Claude models run with adaptive thinking at the Messages API's default effort (Haiku medium, Sonnet high), GPT with medium reasoning effort. Each configuration was run once, with cached answers, apart from the UMR repetitions described above.

\paragraph{Baselines and learned harnesses}
We compare three annotation approaches. The baseline \emph{guideline prompt} makes one call per sentence with the whole guideline and the sentence, using Haiku, Sonnet or GPT, without earlier sentences or examples. The \emph{seed harness $H_0$} uses the three steps written from the guideline, with all four Haiku/GPT generator--critic pairs and with Sonnet in both roles. \emph{\sys} runs ten rounds of RSI with Haiku in both roles, evaluating $H_0,\ldots,H_{10}$. We then transfer $H_{10}$ unchanged to the other generator--critic pairs and to Sonnet; pairwise transfer comparisons use the same two documents per benchmark. For TimeML, the main trajectory trains on platinum documents outside the test set and a second trajectory on TBAQ-cleaned. All approaches receive the same output-format instructions.

\paragraph{Guideline access and harness ablations}
To compare ways of applying the guideline, the \emph{guideline-through-tools} variant supplies only its address and lets the model search and read it, together with UMR resource files; Haiku runs on all test documents and GPT on two per benchmark. To isolate instruction refinement, a \emph{guideline-only round} starts from the guideline prompt and permits only edits or rewrites, with no steps, tools, rules, or length limit. To isolate harness components, we run two \emph{restricted RSI} variants for two rounds each: one disables tool operations, the other rule steps. For TimeML we also test an explicit linking step (\S\ref{sec:link}).

\paragraph{Prior systems}
We re-score earlier systems on exactly our test documents from their released outputs: the two UMR shared-task systems \citep{heinecke2026orange,devergnette2026sema}, fine-tuned on the shared task's training data, of which Orange is best on both Chinese and English; and the TempEval-3 end-to-end systems \citep{uzzaman2013semeval}, of which ClearTK-2 \citep{bethard2013cleartk} is best on temporal awareness and is the best earlier system we compare with. The organizers also report TIPSem \citep{llorens2010tipsem} as a reference that scores higher, and they separate it from the participants because some annotators used it to pre-label the platinum data \citep[\S2.2, \S5]{uzzaman2013semeval}; on our documents 73\% of the gold TLINK pairs are pairs that TIPSem links (ClearTK-2: 26\%). We report TIPSem as a reference, not as the state of the art.
